% !TeX root = ../main.tex
% !TeX program = pdflatex
% !TeX encoding = UTF-8

\chapter{Tập dữ liệu}
\label{chap:dataset}

\section{Nguồn và phương pháp tạo dữ liệu}

Danh sách đầu vào gồm 25 địa điểm du lịch và đầu mối giao thông tại Đà Lạt,
được lưu trong \texttt{data/nodes.geojson}. Script \texttt{build\_edges.py}
dùng OSMnx để tải mạng đường lái xe từ OpenStreetMap, ghép mỗi địa điểm vào
đỉnh đường gần nhất và tính tuyến giữa các cặp địa điểm được chọn. Một cây
khung nhỏ nhất bảo đảm toàn bộ mạng liên thông; mỗi địa điểm còn được nối với
ba láng giềng địa lý gần nhất để tạo các tuyến thay thế.

Khoảng cách và loại đường được lấy từ mạng OpenStreetMap. Thời gian cơ sở được
tính từ khoảng cách và vận tốc mô hình theo loại đường (15--50 km/h), không
phải dữ liệu giao thông trực tiếp. Script \texttt{generate\_edge\_conditions.py}
dùng seed 2026 để sinh điều kiện có thể tái lập cho năm kịch bản S0--S4. Vì vậy
các mức ùn tắc và rủi ro là dữ liệu mô phỏng, không phải quan sát thời gian thực.

Sau tiền xử lý, backend dùng 25 đỉnh, 114 cạnh có hướng và 570 bản ghi điều
kiện (114 cạnh nhân với 5 kịch bản). Các file CSV và GeoJSON sinh ra dùng chung
\texttt{node\_id} và \texttt{edge\_id} để backend và frontend đồng bộ.

\section{Danh sách địa điểm}

\begin{longtable}{p{1.2cm}p{7.0cm}p{4.6cm}}
\caption{Các địa điểm trong tập dữ liệu}\label{tab:locations}\\
\toprule
ID & Tên địa điểm & Loại \\
\midrule
\endfirsthead
\toprule
ID & Tên địa điểm & Loại \\
\midrule
\endhead
DL01 & Nhà lao thiếu nhi Đà Lạt & Bảo tàng \\
DL02 & Ga Đà Lạt & Di sản \\
DL03 & Crazy House & Kiến trúc \\
DL04 & Museum Madam De Dalat & Bảo tàng \\
DL05 & Bảo tàng Lâm Đồng & Bảo tàng \\
DL06 & Thung Lũng Tình Yêu & Khu du lịch \\
DL07 & Thác Datanla & Thác nước \\
DL08 & Devil Mountain Tourist Area & Khu du lịch \\
DL09 & Làng Cu Lần & Khu du lịch \\
DL10 & Vườn dâu Công nghệ cao Bình Yên & Vườn \\
DL11 & Thác Prenn & Thác nước \\
DL12 & Đường hầm điêu khắc & Điểm tham quan \\
DL13 & Dinh III Bảo Đại & Di sản \\
DL14 & Vườn Hoa Cẩm Tú Cầu & Vườn \\
DL15 & Nhà thờ Cam Ly & Công trình tôn giáo \\
DL16 & Hồ Tuyền Lâm & Hồ \\
DL17 & Khu du lịch sinh thái Cao Nguyên Hoa Đà Lạt & Công viên chủ đề \\
DL18 & Làng Cổ Tích & Điểm tham quan \\
DL19 & Làng Đồng trụ LangBiang & Văn hóa \\
DL20 & Thác Ankroet & Thác nước \\
DL21 & Chợ Đà Lạt & Chợ \\
DL22 & Hồ Xuân Hương & Hồ \\
DL23 & Quảng trường Lâm Viên & Địa danh \\
DL24 & Thiền viện Trúc Lâm Đà Lạt & Công trình tôn giáo \\
DL25 & Bến xe Liên tỉnh Đà Lạt & Giao thông \\
\bottomrule
\end{longtable}

\section{Khoảng cách, thời gian, ùn tắc và loại đường}

Khoảng cách được lưu theo km với ba chữ số thập phân; thời gian cơ sở được lưu
theo phút. Mỗi kịch bản gán mức ùn tắc nguyên từ 1 (thông thoáng) đến 5 (ùn tắc
nặng), hệ số thời gian, rủi ro mưa, sương mù, thi công và trạng thái đóng/mở cho
từng cạnh. Giá trị thay đổi theo kịch bản nhưng cố định trong một lần tìm đường.

\begin{table}[h!]
\centering
\caption{Một số cạnh đầu tiên trong dữ liệu gốc}
\small
\begin{tabular}{p{1.1cm}p{2.6cm}rrp{4.4cm}}
\toprule
Cạnh & Hướng & Khoảng cách (km) & Thời gian (phút) & Loại đường \\
\midrule
E001 & DL01 $\rightarrow$ DL02 & 2,320 & 3,56 & residential, secondary \\
E002 & DL02 $\rightarrow$ DL01 & 2,320 & 3,56 & residential, secondary \\
E003 & DL01 $\rightarrow$ DL05 & 2,207 & 4,24 & residential, secondary, trunk \\
E005 & DL01 $\rightarrow$ DL06 & 6,008 & 9,10 & residential, secondary \\
E007 & DL01 $\rightarrow$ DL14 & 7,289 & 15,20 & primary, residential, secondary, trunk, unclassified \\
\bottomrule
\end{tabular}
\end{table}

\section{Các giả định và đơn giản hóa}

\begin{itemize}[itemsep=2pt]
  \item Tập dữ liệu chỉ nối các địa điểm đã chọn, không chứa mọi giao lộ của thành
  phố như những đỉnh tìm kiếm độc lập.
  \item Vận tốc theo loại đường là giả định mô hình, không phải giới hạn tốc độ
  được công bố hoặc số đo từ thiết bị GPS.
  \item Ùn tắc, thời tiết, thi công và đóng đường được sinh giả lập bằng seed cố
  định; chúng không thay đổi trong thời gian một chuyến đi.
  \item Thời gian chờ đèn giao thông và tại giao lộ không được mô hình riêng.
  \item Hình học cạnh hiện được xuất thành đoạn thẳng giữa hai địa điểm để giao
  diện nhẹ hơn, trong khi khoảng cách và thời gian vẫn dựa trên tuyến đường đã
  tính từ OpenStreetMap.
\end{itemize}
